\chapter{Nowoczesne niezmienniki i zastosowania teorii węzłów}
\label{ch:kt-nowoczesne}

Klasyczne wielomiany i grupy wystarczają do rozwiązania wielu
problemów, ale nie rozróżniają wszystkich subtelności gładkiej
topologii czterowymiarowej. Współczesne niezmienniki przypisują
węzłowi nie tylko liczbę albo wielomian, lecz cały kompleks
łańcuchowy i jego homologie. Dają one nowe ograniczenia
na gładki rodzaj czterowymiarowy oraz rozpoznają własności
takie jak włóknistość.

\section{Homologia Khovanova: od wielomianu do kompleksu}
\label{sec:kt-khovanov}

\begin{definicja}[Algebra Frobeniusa]
Niech $A:=\mathbb Z[X]/(X^2)$ będzie wolną grupą abelową
z bazą $(1,X)$ i gradacją $\deg(1)=1$, $\deg(X)=-1$.
Mnożenie $m\colon A\otimes A\to A$ jest indukowane przez
mnożenie wielomianów, a komnożenie
$\Delta\colon A\to A\otimes A$ jest dane wzorami
$\Delta(1)=1\otimes X+X\otimes1$,
$\Delta(X)=X\otimes X$. Odwzorowania $m$ i $\Delta$
są jednorodne stopnia $-1$ względem powyższej gradacji.
\end{definicja}

\begin{lemat}[Tożsamości algebry stanów Khovanova]
\label{lem:kt-khovanov-frobenius}
Dla $A=\mathbb Z[X]/(X^2)$ z mnożeniem $m$
i komnożeniem $\Delta$ podanymi wyżej zachodzą:
\[
m(m\otimes\mathrm{id})
=m(\mathrm{id}\otimes m),\qquad
(\Delta\otimes\mathrm{id})\Delta
=(\mathrm{id}\otimes\Delta)\Delta,
\]
oraz
\[
\Delta m=(m\otimes\mathrm{id})
(\mathrm{id}\otimes\Delta)
=(\mathrm{id}\otimes m)
(\Delta\otimes\mathrm{id}).
\]
Wszystkie odwzorowania w tych równościach
mają odpowiednie wspólne dziedziny i
przeciwdziedziny: pierwsza para dotyczy
$A^{\otimes3}\to A$ i $A\to A^{\otimes3}$,
a ostatnie odwzorowania działają
z $A^{\otimes2}$ do $A^{\otimes2}$.
Są to algebraiczne przesłanki przemienności
dwuwymiarowych ścian kostki Khovanova.
\end{lemat}

\begin{proof}
Łączność $m$ wynika z łączności pierścienia
$\mathbb Z[X]/(X^2)$. Współłączność $\Delta$
wystarczy sprawdzić na bazie $(1,X)$.
Dla $1$ obydwa złożenia dają
\[
1\otimes X\otimes X
+X\otimes1\otimes X
+X\otimes X\otimes1,
\]
zaś dla $X$ otrzymujemy
$X\otimes X\otimes X$.

Tożsamość Frobeniusa sprawdzamy na
czterech elementach bazy
$(1\otimes1,1\otimes X,
X\otimes1,X\otimes X)$.
Trzy wyrażenia z tezy mają
odpowiednio wspólne wartości
\[
1\otimes X+X\otimes1,\qquad
X\otimes X,\qquad
X\otimes X,\qquad 0.
\]
Rachunek wykorzystuje tylko
$m(1\otimes1)=1$,
$m(1\otimes X)=m(X\otimes1)=X$,
$m(X\otimes X)=0$ oraz
dwa wzory definiujące $\Delta$.
Ponieważ rozważane odwzorowania
są $\mathbb Z$-liniowe,
sprawdzenie na bazie
kończy dowód.
\end{proof}

\begin{definicja}[Stany kostki rozwiązań]
Niech $D$ będzie zorientowanym diagramem splotu
$L\subset S^3$ o $n$ uporządkowanych skrzyżowaniach.
Dla każdego $s=(s_1,\dots,s_n)\in\{0,1\}^n$
wygładzamy skrzyżowanie numer $j$ według konwencji
0 lub 1: wygładzenie 0 jest zgodne z orientacją
pasm przy dodatnim skrzyżowaniu, a wygładzenie 1
jest zgodne z orientacją przy ujemnym skrzyżowaniu.
Drugie wygładzenie w każdym przypadku jest tym
pozostałym sposobem połączenia czterech końców. Otrzymujemy rozłączną sumę $k(s)$ okręgów.
Niech $|s|:=s_1+\dots+s_n$.
Przypisujemy stanowi moduł $A^{\otimes k(s)}$, przyjmując
$A^{\otimes0}=\mathbb Z$. Oznaczmy przez $n_+$ oraz $n_-$
liczby dodatnich i ujemnych skrzyżowań diagramu.
\end{definicja}

Dla jasności podajemy jedną z najczęściej stosowanych
konwencji przesunięć gradacji. Jeżeli $M\{r\}$ oznacza
moduł, w którym stopnie elementów zostały zwiększone o $r$,
to definiujemy grupę kompleksu w stopniu kohomologicznym $i$
(tradycyjnie w teorii Khovanova nazywanym
stopniem homologicznym):
\[
C^{i}(D):=
\bigoplus_{\substack{s\in\{0,1\}^n\\|s|-n_-=i}}
A^{\otimes k(s)}\{|s|+n_+-2n_-\}.
\]
Krawędź kostki zmieniająca jedno 0 w 1 odpowiada
albo połączeniu dwóch okręgów w jeden, albo rozszczepieniu
jednego okręgu na dwa. Odpowiedni homomorfizm jest więc
$m$ albo $\Delta$, z identycznościami na pozostałych
czynnikach. W każdym stanie ustalamy kolejność
okręgów, a przy stosowaniu $m$ lub $\Delta$
odpowiednio permutujemy czynniki tensorowe
(bez znaków Koszula: gradacja kwantowa jest
niezależna od znaku różniczki).
Sumując te homomorfizmy ze znakami
kostki, otrzymujemy różniczkę $d\colon C^i(D)\to C^{i+1}(D)$.

\begin{stwierdzenie}[Dwugradacja różniczki i charakterystyka Eulera]
\label{prop:kt-kh-graded-euler}
Niech $D$ będzie zorientowanym diagramem splotu
o uporządkowanych $n$ skrzyżowaniach,
z $n_+$ skrzyżowaniami dodatnimi i
$n_-$ ujemnymi. Zdefiniujmy
dwugradowany kompleks $C^{i,j}(D)$
zgodnie z poprzednią konstrukcją.
Wtedy każda krawędź różniczki ma
bistopień $(1,0)$, czyli
$d:C^{i,j}(D)\to C^{i+1,j}(D)$.
Jeżeli $q$ jest zmienną formalną, to
gradująca charakterystyka Eulera
homologii Khovanova ma postać
\[
\sum_{i,j}(-1)^i q^j
\operatorname{rank}_{\mathbb Z}
\operatorname{Kh}^{i,j}(D;\mathbb Z)
=
\sum_{s\in\{0,1\}^n}
(-1)^{|s|-n_-}
q^{|s|+n_+-2n_-}
(q+q^{-1})^{k(s)}.
\]
W szczególności dla diagramu
węzła trywialnego bez skrzyżowań
otrzymujemy $q+q^{-1}$, a dla
dwuskładowego splotu trywialnego
$(q+q^{-1})^2$.
\end{stwierdzenie}

\begin{proof}
Homomorfizmy $m$ i $\Delta$ mają
stopień kwantowy $-1$
przed uwzględnieniem przesunięć.
Wzdłuż krawędzi kostki liczba
$|s|$ wzrasta o jeden, dlatego
przesunięcie gradacji na
przeciwdziedzinie zwiększa
stopień kwantowy o jeden.
Łączny stopień kwantowy
wynosi $-1+1=0$,
a stopień homologiczny wzrasta
o jeden. Tak otrzymujemy
bistopień $(1,0)$.

Dla modułu $A$ mamy
$\operatorname{grdim}_q A=q+q^{-1}$.
Zatem stan z $k(s)$ okręgami
wnosi do charakterystyki
Eulera kompleksu składnik
$q^{|s|+n_+-2n_-}(q+q^{-1})^{k(s)}$
ze znakiem homologicznym
$(-1)^{|s|-n_-}$.
W każdym ustalonym stopniu
kwantowym kompleks jest
skończonym kompleksem
skończenie generowanych
wolnych grup abelowych.
Po przejściu do $\mathbb Q$
(funktor $-\otimes_{\mathbb Z}\mathbb Q$ jest dokładny,
więc homologia kompleksu nad $\mathbb Q$ jest
iloczynem tensorowym homologii całkowitej z $\mathbb Q$)
zasada Eulera--Poincarégo
utożsamia naprzemienną
sumę rang grup kompleksu
z naprzemienną sumą
rang homologii. To dowodzi
wzoru bez założenia zanikania
możliwej torsji homologii.

Dla diagramu bez skrzyżowań
jedyny stan ma $|s|=0$.
Liczba okręgów wynosi
odpowiednio $1$ lub $2$,
co daje oba podane przykłady.
\end{proof}

\begin{przyklad}[Charakterystyka Eulera dla dodatniego i ujemnego skrętu]
\label{ex:kt-kh-reidemeister-one-euler}
Niech $U$ będzie diagramem węzła trywialnego
bez skrzyżowań i niech $D_+$, $D_-$ będą
diagramami otrzymanymi przez wykonanie
odpowiednio jednego dodatniego i
jednego ujemnego ruchu Reidemeistera typu I.
Oznaczmy $f(q):=q+q^{-1}$.
Przy konwencji wygładzeń z tego rozdziału
dodatnie skrzyżowanie ma w stanie $0$
dwa okręgi, a w stanie $1$ jeden okrąg.
Dlatego wzór ze
Stwierdzenia~\ref{prop:kt-kh-graded-euler}
daje
\[
\chi_q(C(D_+))=qf(q)^2-q^2f(q)=f(q).
\]
Przy ujemnym skrzyżowaniu stan $0$
ma jeden okrąg, a stan $1$ dwa.
Ponieważ wtedy $n_+=0$, $n_-=1$,
dostajemy analogicznie
\[
\chi_q(C(D_-))
=-q^{-2}f(q)+q^{-1}f(q)^2=f(q).
\]
Obliczenia sprawdzają normalizację
gradującej charakterystyki Eulera
dla obu skrętów. Nie są jeszcze
dowodem niezmienniczości \emph{kompleksu}
Khovanova względem ruchów Reidemeistera:
ten silniejszy wynik wymaga
porównania różniczek i równoważności
homotopijnych.
\end{przyklad}

\begin{przyklad}[Homologia obu skrętów węzła trywialnego]
\label{ex:kt-kh-reidemeister-one-homology}
Dla diagramów $D_+$ i $D_-$ z
Przykładu~\ref{ex:kt-kh-reidemeister-one-euler}
można sprawdzić nie tylko charakterystyki Eulera,
lecz również obie gradacje homologii.
Dla $D_+$ kompleks ma postać
\[
 A^{\otimes2}\{1\}\xrightarrow{m}A\{2\}
\]
w stopniach $0,1$. Mnożenie jest surjektywne,
a jego jądro jest wolne z bazą
$1\otimes X-X\otimes1,\ X\otimes X$:
dla elementu o współczynnikach $(a,b,c,d)$
w standardowej bazie iloczynu tensorowego obraz
wynosi $a+(b+c)X$, więc $a=0$ i $b=-c$.
Po przesunięciu stopnie kwantowe generatorów
jądra wynoszą odpowiednio $1,-1$.
Zatem homologia występuje tylko w stopniu $0$.

Dla $D_-$ kompleks ma postać
\[
 A\{-2\}\xrightarrow{\Delta}A^{\otimes2}\{-1\}
\]
w stopniach $-1,0$. Obrazy $1,X$ są liniowo
niezależne, a iloraz przez obraz $\Delta$
jest wolny z bazą klas $1\otimes1,\ 1\otimes X$:
relacje to $X\otimes X=0$ oraz
$X\otimes1=-1\otimes X$, bez dodatkowych relacji
między wybranymi generatorami. Ich stopnie kwantowe
po przesunięciu wynoszą $1,-1$.
W obu przypadkach otrzymujemy zatem
$\operatorname{Kh}^{0,1}\cong
\operatorname{Kh}^{0,-1}\cong\mathbb Z$,
a pozostałe grupy zanikają.
Rachunek dotyczy obu skrętów diagramu węzła
trywialnego i nie zastępuje dowodu ogólnej
niezmienniczości kompleksu względem ruchów Reidemeistera.
\end{przyklad}

\begin{uwaga}[Znaki różniczki i jej obliczanie]
Uporządkujmy skrzyżowania $1,\dots,n$. Jeżeli stan
$s\in\{0,1\}^n$ ma $s_j=0$, to krawędzi zmieniającej
$s_j$ na $1$ przypisujemy znak
$(-1)^{\sum_{k<j}s_k}$. Pozostała część operatora tej
krawędzi jest odpowiednim mnożeniem $m$ albo
komnożeniem $\Delta$. Dla dowolnych $j<k$ dwie
drogi w kwadracie kostki otrzymują przeciwne znaki,
ponieważ w jednej kolejności zmiana $s_j$ zwiększa
o jeden wykładnik znaku przy zmianie $s_k$, a w drugiej
nie. Otrzymujemy w ten sposób jednoznaczną, skończoną
procedurę zapisu macierzy różniczki.
\end{uwaga}

\begin{stwierdzenie}[Różniczka kwadratuje się do zera]
Dla dowolnego zorientowanego diagramu splotu $D$ i ustalonego
porządku jego skrzyżowań operator kostkowy
$d:C^i(D)\to C^{i+1}(D)$ spełnia $d^2=0$.
\end{stwierdzenie}

\begin{proof}
Każda składowa $d^2$ prowadząca ze stanu $s$ do stanu,
w którym dwa różne skrzyżowania $j,k$ zmieniły wartość
z $0$ na $1$, jest sumą dwóch kompozycji: najpierw $j$,
potem $k$, oraz w odwrotnej kolejności. Operacje
łączenia i rozszczepiania okręgów opisane przez $m$
i $\Delta$ tworzą przemienny kwadrat algebraiczny.
Jeśli oba siodła działają na rozłącznych
zbiorach okręgów, odwzorowania komutują,
gdyż działają na niezależnych czynnikach
iloczynu tensorowego. Dla połączonych siodeł
dwa łączenia dają łączność $m$, dwa rozdzielenia
współłączność $\Delta$, a układ z jednym
łączeniem i jednym rozdzieleniem przy dwóch
okręgach wejściowych i dwóch wyjściowych
daje tożsamość Frobeniusa
\[
\Delta m=(m\otimes\mathrm{id})(\mathrm{id}\otimes\Delta)
=(\mathrm{id}\otimes m)(\Delta\otimes\mathrm{id}).
\]
Pozostały połączony przypadek, w którym
na wejściu i wyjściu jest po jednym okręgu,
daje w obu kolejnościach $m\Delta$.
Permutacje wejść i wyjść nie zmieniają tych
wniosków: $m$ jest przemienne, a $\Delta$
współprzemienne, co widać z podanych wzorów.
Przy $j<k$ wykładniki znaku przy drogach $j\to k$
i $k\to j$ różnią się dokładnie o jeden: zmiana $s_j$
powoduje dodatkowy składnik w sumie $\sum_{\ell<k}s_\ell$,
a zmiana $s_k$ nie wpływa na znak krawędzi $j$.
Dlatego obie kompozycje mają przeciwne znaki
i ich suma jest zerowa. Nie ma składnika $d^2$
zmieniającego tylko jeden stan, ponieważ każda krawędź
kostki zwiększa liczbę jedynek o jeden. Stąd $d^2=0$.
\end{proof}

\begin{twierdzenie}[Niezmienniczość homologii Khovanova; wynik zewnętrzny]
Po uwzględnieniu wskazanych przesunięć gradacji
homologia kompleksu $(C^*(D),d)$ nie zależy od wyboru
regularnego zorientowanego diagramu danego splotu,
z dokładnością do izomorfizmu zachowującego obie gradacje.
\end{twierdzenie}

\begin{uwaga}[Zakres dowodu]
Pełna niezmienniczość wymaga skonstruowania dla każdego
z trzech ruchów Reidemeistera równoważności homotopijnych
odpowiednich kompleksów, z rozdzieleniem składników
kontraktowalnych i kontrolą przesunięć gradacji.
Nie jest wnioskiem z samego równania $d^2=0$.
Przyjmujemy ją w tym miejscu jako twierdzenie Khovanova;
poprzedni dowód dotyczy tylko formalnej konstrukcji kompleksu.
\end{uwaga}

\begin{definicja}[Homologia Khovanova]
Dla zorientowanego splotu $L\subset S^3$ przez
$\operatorname{Kh}^{i,j}(L;\mathbb Z)$ oznaczamy dwugradowaną
homologię kompleksu zbudowanego z dowolnego diagramu $L$
według poprzedniej konstrukcji, w stopniach homologicznym $i$
i kwantowym $j$. \emph{Zredukowana homologia Khovanova}
$\widetilde{\operatorname{Kh}}(L,p;\mathbb Z)$
wymaga wyboru punktu bazowego $p$ na
jednej składowej diagramu. W każdym stanie
wybieramy czynnik $A$ odpowiadający
okręgowi zawierającemu $p$ i ograniczamy
go do ideału $XA\subset A$,
z przesunięciem gradacji kwantowej o $1$.
Daje to podkompleks
$\widetilde C(D,p)\subset C(D)\{1\}$.
Jest on zachowywany przez różniczkę:
przy sklejeniu okręgów $m(X\otimes A)\subset XA$,
a przy rozdzieleniu oznaczonego okręgu
$\Delta(X)=X\otimes X$,
więc czynnik zawierający $p$
nadal leży w $XA$.
Dla węzła otrzymana dwugradowana
homologia nie zależy od wyboru
punktu bazowego (do izomorfizmu);
dla splotu wybór składowej bazowej
może wpływać na wersję zredukowaną.
Niezmienniczość tej konstrukcji
wobec ruchów Reidemeistera
przyjmujemy zewnętrznie,
tak jak dla kompleksu niezredukowanego.
\end{definicja}

\begin{stwierdzenie}[Dokładna normalizacja wielomianu Jonesa]
\label{prop:kt-kh-jones-normalization}
Niech $L$ będzie niepustym zorientowanym splotem.
W konwencji wielomianu Jonesa z Sekcji~\ref{sec:kt-jones}
oznaczmy $J_L(q):=V_L(t)|_{t^{1/2}=-q}$, czyli
podstawiamy $t^{m/2}=(-q)^m$ dla $m\in\mathbb Z$.
Dla węzłów sprowadza się to do podstawienia $t=q^2$.
Wtedy
\[
 \chi_q(\operatorname{Kh}(L;\mathbb Z))
 =(q+q^{-1})J_L(q),\qquad
 \chi_q(\widetilde{\operatorname{Kh}}(L,p;\mathbb Z))
 =J_L(q).
\]
Przez $\chi_q$ rozumiemy naprzemienną sumę rang
w stopniach homologicznym i kwantowym, jak w
Stwierdzeniu~\ref{prop:kt-kh-graded-euler}.
\end{stwierdzenie}

\begin{proof}
Użyjmy zmiennej $a$ dla nawiasu Kauffmana, aby
odróżnić ją od algebry stanów $A$.
Przy przyjętej konwencji wygładzenie $0$ jest
wygładzeniem typu A, a wygładzenie $1$ typu B
w Sekcji~\ref{sec:kt-jones}.
Dla $r=|s|$, $k=k(s)$, $n=n_++n_-$,
oraz $w=n_+-n_-$, rozwinięcie nawiasu daje
\[
 V_L(t)\big|_{t^{1/2}=a^{-2}}
 =(-a^3)^{-w}
 \sum_s a^{n-2r}(-a^2-a^{-2})^{k-1}.
\]
Podstawienie $q=-a^{-2}$ daje
$q+q^{-1}=-a^2-a^{-2}$.
Po pomnożeniu powyższej równości przez $q+q^{-1}$
współczynnik przy $(q+q^{-1})^k$ wynosi
\[
 (-1)^w a^{-3w+n-2r}
 =(-1)^w a^{-2r-2n_++4n_-}.
\]
Z drugiej strony
\[
 \begin{aligned}
 (-1)^{r-n_-}q^{r+n_+-2n_-}
 &=(-1)^{2r+n_+-3n_-}
   a^{-2r-2n_++4n_-}\\
 &=(-1)^w a^{-2r-2n_++4n_-},
 \end{aligned}
\]
gdyż $2r+n_+-3n_-\equiv w\pmod 2$.
Porównanie składników z
\ref{prop:kt-kh-graded-euler} dowodzi pierwszej
równości po podstawieniu $q=-a^{-2}$, które
jest różnowartościowe na pierścieniu $\mathbb Z[q,q^{-1}]$.
W wersji zredukowanej oznaczony czynnik $XA\{1\}$
ma jedyny generator stopnia zero.
W każdym stanie $k\geq1$ i czynnik
$(q+q^{-1})^k$ zastępuje się więc przez
$(q+q^{-1})^{k-1}$, bez zmiany znaku i pozostałych
przesunięć. Ta sama zasada Eulera--Poincarégo
dowodzi drugiej równości.
\end{proof}

\begin{uwaga}[Kategoryfikacja]

Dwa sploty mogą mieć jednakowy wielomian Jonesa, lecz
różne homologie Khovanova (por. D.~Bar-Natan,
\emph{On Khovanov's categorification of the Jones polynomial},
Algebr. Geom. Topol. 2 (2002), 337--370).
Homologia zachowuje informacje
pominięte przy obliczeniu samej charakterystyki Eulera.
\end{uwaga}

\begin{przyklad}[Węzeł trywialny]
Diagram węzła trywialnego nie ma skrzyżowań.
Jego kompleks Khovanova składa się wyłącznie z
$C^0(U)=A$ i zerowej różniczki. Stąd
$\operatorname{Kh}(U;\mathbb Z)\cong\mathbb Z\oplus\mathbb Z$
z generatorami w stopniach kwantowych $\pm1$
i stopniu homologicznym zero.
Po redukcji jedyny generator $X$ ma stopień
kwantowy $-1+1=0$ i stopień homologiczny zero,
więc $\widetilde{\operatorname{Kh}}(U;\mathbb Z)
\cong\mathbb Z$ w bistopniu $(0,0)$.
\end{przyklad}

\begin{twierdzenie}[Kronheimer--Mrowka; homologia Khovanova wykrywa węzeł trywialny]
Niech $K\subset S^3$ będzie węzłem. Jeżeli
$\operatorname{rank}_{\mathbb Q}
\widetilde{\operatorname{Kh}}(K;\mathbb Q)=1$,
to $K$ jest trywialny. Odwrotność wynika bezpośrednio
z obliczenia homologii Khovanova dla węzła trywialnego.
\end{twierdzenie}

\begin{uwaga}
To twierdzenie jest zdecydowanie silniejsze niż
stwierdzenie, że odpowiednio znormalizowany wielomian
Jonesa węzła trywialnego wynosi jeden. Nie wiadomo,
czy warunek $V_K(t)=1$ sam w sobie wystarcza
rozstrzygnąć trywialność $K$; homologia Khovanova
rozpoznaje ją na poziomie dodatkowej struktury homologicznej.
Dowód Kronheimera--Mrowki wykorzystuje teorię instantonów.
\end{uwaga}

\section{Niezmiennik Rasmussena i gładki rodzaj}
\label{sec:kt-rasmussen}

Deformację Lee konstruujemy nad $\mathbb Q$, zachowując
bazę $(1,X)$ i grupy kompleksu kostkowego, lecz zmieniając
algebrę na $\mathbb Q[X]/(X^2-1)$. Nowe mnożenie
$m_L$ i komnożenie $\Delta_L$ różnią się od poprzednich
tylko wartościami
\[
 m_L(X\otimes X)=1,\qquad
 \Delta_L(X)=X\otimes X+1\otimes1.
\]
Spełniają te same tożsamości łączności,
współłączności i Frobeniusa co poprzednie
odwzorowania. Na przykład na bazie
$(1\otimes1,1\otimes X,X\otimes1,X\otimes X)$
trzy strony tożsamości Frobeniusa mają wartości
$1\otimes X+X\otimes1$, $X\otimes X+1\otimes1$,
$X\otimes X+1\otimes1$, $1\otimes X+X\otimes1$.
Dla współłączności rozwijamy obie kompozycje
$(\Delta_L\otimes\mathrm{id})\Delta_L$ oraz
$(\mathrm{id}\otimes\Delta_L)\Delta_L$ na bazie:
\[
\begin{aligned}
1&\longmapsto 1\otimes X\otimes X+
 X\otimes1\otimes X+X\otimes X\otimes1+
 1\otimes1\otimes1,\\
X&\longmapsto X\otimes X\otimes X+
 1\otimes1\otimes X+1\otimes X\otimes1+
 X\otimes1\otimes1.
\end{aligned}
\]
Łączność pochodzi z mnożenia w pierścieniu, zaś
przemienność mnożenia i współprzemienność komnożenia
wynikają bezpośrednio z podanych wzorów.
Dowód z kwadratów kostki daje więc $d_L^2=0$.

\begin{uwaga}[Baza idempotentów algebry Lee]
\label{rem:kt-lee-idempotents}
W algebrze $\mathbb Q[X]/(X^2-1)$ połóżmy
$e_+=(1+X)/2$ oraz $e_-=(1-X)/2$.
Z relacji $X^2=1$ wynika
$e_+^2=e_+$, $e_-^2=e_-$, $e_+e_-=0$
i $e_++e_-=1$. Ponieważ $X=e_+-e_-$,
jest to baza nad $\mathbb Q$. Bezpośrednie podstawienie
do wzorów na komnożenie daje
\[
 \Delta_L(e_+)=2e_+\otimes e_+,\qquad
 \Delta_L(e_-)=-2e_-\otimes e_-.
\]
Mnożenie dwóch różnych idempotentów jest zerowe,
a mnożenie dwóch jednakowych daje ten sam idempotent.
Współczynniki $\pm2$ są odwracalne nad $\mathbb Q$,
co jest użyteczne przy rozkładaniu kompleksu Lee.
Nie jest to jednak baza jednorodna względem gradacji
kwantowej: każdy z tych elementów zawiera składniki
stopni $1$ i $-1$. Sama zmiana bazy nie dowodzi
twierdzenia Lee o wymiarze homologii.
\end{uwaga}

Dodatkowe składniki $m_L,\Delta_L$ mają stopień
kwantowy $3$ przed przesunięciem, a zatem stopień
$4$ po przesunięciu. Różniczka $d_L$ jest sumą
różniczki Khovanova stopnia kwantowego zero
i składnika stopnia cztery.
Zapiszmy $d_L=d+d_4$, gdzie $d$ jest różniczką
Khovanova, a $d_4$ podnosi stopień kwantowy o cztery.
Składowe równania $d_L^2=0$ w stopniach kwantowych
$0,4,8$ muszą znikać oddzielnie, ponieważ kompleks
jest sumą prostą grup o różnych stopniach. Otrzymujemy
\[
 d^2=0,\qquad dd_4+d_4d=0,\qquad d_4^2=0.
\]
Różniczka $d_L$ nie zachowuje na ogół gradacji kwantowej, lecz
zachowuje malejącą filtrację
$F^j C_L=\bigoplus_{k\ge j} C^{*,k}(D;\mathbb Q)$.
Na homologii przyjmujemy filtrację
$F^j H_L=\operatorname{im}(H^*(F^j C_L)\to H^*(C_L))$.
Z twierdzenia Lee homologia węzła nad $\mathbb Q$
ma wymiar dwa. Jeżeli $v\ne0$ jest klasą w tej
homologii, kładziemy
$\ell(v)=\max\{j:v\in F^j H_L\}$.
Minimum i maksimum wartości $\ell(v)$ dla $v\ne0$
oznaczamy przez $s_{\min}$ i $s_{\max}$; istnieją,
bo filtracja jest skończona i wyczerpująca.
Rasmussen wykazał, że te liczby są niezmiennikami
węzła, są nieparzyste i $s_{\max}-s_{\min}=2$.
Definiujemy
\[
 s(K)=\frac{s_{\min}(K)+s_{\max}(K)}{2}\in2\mathbb Z.
\]
W przypadku węzła trywialnego różniczka jest zerowa,
a dwa poziomy filtracji to $-1,1$, więc $s(U)=0$.
Wyniki o wymiarze, niezmienniczości i różnicy
poziomów przyjmujemy z pracy J.~Rasmussena,
\emph{Khovanov homology and the slice genus},
Invent. Math. 182 (2010), 419--447, sekcje 2--3
(arXiv:math/0402131).

\begin{twierdzenie}[Rasmussen]
Zdefiniowany powyżej niezmiennik $s$ spełnia
następujące własności (przyjmujemy twierdzenia
1, 2 i 4 oraz własności z sekcji 3 cytowanej
pracy Rasmussena). Przez $\overline K$ rozumiemy
odbicie lustrzane węzła, a $g_4^{\mathrm{sm}}(K)$
oznacza najmniejszy rodzaj spójnej zorientowanej
gładko właściwie zanurzonej powierzchni w $D^4$
o brzegu $K$: 
\begin{enumerate}[label=(\roman*)]
\item $s(U)=0$;
\item $s(K_1\#K_2)=s(K_1)+s(K_2)$;
\item $s(\overline K)=-s(K)$;
\item $|s(K)|\leq2g_4^{\mathrm{sm}}(K)$;
\item dla dodatnich względnie pierwszych $p,q\geq2$,
$s(T(p,q))=(p-1)(q-1)$.
\end{enumerate}
\end{twierdzenie}

\begin{wniosek}
Dla dodatnich węzłów torusowych $T(p,q)$,
$p,q\geq2$ względnie pierwszych, mamy
$g_4^{\mathrm{sm}}(T(p,q))=(p-1)(q-1)/2$.
\end{wniosek}

\begin{proof}
Z nierówności Rasmussena i jego obliczenia dla węzłów
torusowych otrzymujemy
$(p-1)(q-1)/2=|s(T(p,q))|/2\leq g_4^{\mathrm{sm}}(T(p,q))$.
Z drugiej strony $g_4^{\mathrm{sm}}(K)\leq g(K)$:
powierzchnię Seiferta można przesunąć we wnętrzu
kołnierza $S^3=\partial D^4$, pozostawiając jej
brzeg na miejscu. Rodzaj powierzchni się nie zmienia.
Dla górnego ograniczenia nie potrzeba niezależnego
twierdzenia o minimalności rodzaju powierzchni Seiferta.
Węzeł $T(p,q)$ jest domknięciem dodatniego warkocza
$(\sigma_1\cdots\sigma_{p-1})^q$. Algorytm Seiferta
daje $p$ dysków (po jednym dla każdego okręgu
zorientowanego wygładzenia) oraz $q(p-1)$ wstęg
odpowiadających skrzyżowaniom. Powierzchnia jest
spójna, bo wstęgi łączą kolejne dyski $i,i+1$
dla każdego $1\leq i<p$. Jest zorientowana
i ma jeden składnik brzegu, ponieważ $p,q$ są
względnie pierwsze. Jej charakterystyka Eulera
wynosi $p-q(p-1)$: każdy dysk wnosi $1$,
a dołączenie każdej wstęgi jako uchwytu
indeksu jeden zmniejsza charakterystykę o $1$.
Z równania $1-2g=p-q(p-1)$ wynika, że rodzaj
tej powierzchni wynosi $(p-1)(q-1)/2$.
Stąd $g_4^{\mathrm{sm}}(T(p,q))\leq
g(T(p,q))\leq(p-1)(q-1)/2$.
W połączeniu z dolnym ograniczeniem Rasmussena
wszystkie trzy wielkości są równe.
\end{proof}

\begin{uwaga}
Niezmiennik $s$ jest przykładem rezultatu typowo
czterowymiarowego wydobytego z teorii diagramowej.
Jego obliczanie wymaga homologii i filtracji, ale nie
wymaga bezpośredniej analizy równań teorii cechowania.
Jest to przykład kombinatorycznego dowodu
ograniczenia czterowymiarowego, choć twierdzenia
o filtracji i odwzorowaniach kobordyzmów pozostają
nietrywialną częścią teorii.
\end{uwaga}

\section{Homologia Floera węzła}
\label{sec:kt-hfk}

\begin{definicja}[Kompleks Floera węzła; zakres konstrukcji]
Niech $K\subset S^3$ będzie zorientowanym węzłem, a
$\mathbb F=\mathbb Z/2\mathbb Z$. Wybieramy dopuszczalny
dwupunktowy diagram Heegaarda $(\Sigma,\boldsymbol\alpha,
\boldsymbol\beta,w,z)$ opisujący parę $(S^3,K)$.
Jego generatory są punktami przecięcia torusów
$\mathbb T_{\alpha}$ i $\mathbb T_{\beta}$ w
$\operatorname{Sym}^{g}(\Sigma)$, gdzie $g$ jest
rodzajem powierzchni Heegaarda.
Na wolnym module nad $\mathbb F[U]$ z tymi
generatorami definiujemy filtrowaną różniczkę
\[
 \partial x=\sum_y
 \sum_{\substack{\phi\in\pi_2(x,y)\\\mu(\phi)=1}}
 \#_2\widehat{\mathcal M}(\phi)\,U^{n_w(\phi)}y.
\]
Symbol $\#_2$ oznacza liczbę punktów
zerowymiarowej przestrzeni moduli dysków
Whitneya modulo dwa (po podzieleniu przez
translacje). Liczby $n_w(\phi),n_z(\phi)$ są
krotnościami przecięć z dzielnikami
wyznaczonymi przez punkty bazowe. Dla dysków
pseudoholomorficznych są nieujemne.
Relatywną gradację Alexandera ustala wzór
$A(x)-A(y)=n_z(\phi)-n_w(\phi)$;
normalizację absolutną przyjmujemy standardową
dla węzłów w $S^3$. Stopień Maslova spełnia
$M(x)-M(y)=\mu(\phi)-2n_w(\phi)$,
a $M(U^kx)=M(x)-2k$. Zatem każdy składnik
różniczki ma stopień Maslova o jeden mniejszy
od stopnia wejścia. Definiujemy
$A(U^k x)=A(x)-k$. W konsekwencji
$A(U^{n_w(\phi)}y)=A(x)-n_z(\phi)$,
więc różniczka nie zwiększa filtracji.
Otrzymujemy filtrowany kompleks $CF^-(S^3)$
z filtracją wyznaczoną przez $K$.

W stowarzyszonym kompleksie stopniowanym tej
filtracji pozostają dokładnie składniki
z $n_z(\phi)=0$. Ten kompleks oznaczamy
przez $CFK^-(K)$, a jego homologię przez $HFK^-(K)$.
Nie wolno utożsamiać jego różniczki z pełną
różniczką filtrowaną.

Po przejściu do ilorazu pełnego kompleksu przez
$U$ otrzymujemy filtrowany kompleks
$\widehat{CF}(S^3)$, w którego różniczce pozostają
składniki z $n_w(\phi)=0$. W jego stowarzyszonym
kompleksie stopniowanym pozostają składniki
spełniające jednocześnie $n_w(\phi)=n_z(\phi)=0$.
Przez $\widehat{HFK}(K,s)$ oznaczamy homologię
części tego kompleksu w stopniu Alexandera $s$.
Rozróżnienie między filtrowanym kompleksem
i stowarzyszonym kompleksem stopniowanym jest
istotne dla definicji $\tau$.
Por. P.~Ozsváth, Z.~Szabó, \emph{Knot Floer homology
and rational surgeries}, Algebr. Geom. Topol. 11
(2011), 1--68, sekcja 3. Klasa izomorfizmu grup $\widehat{HFK}(K,s)$
nie zależy od wyboru dopuszczalnego diagramu Heegaarda.
\end{definicja}

\begin{uwaga}
Powyższa konstrukcja nie jest elementarna: do
pełnego dowodu, że $\partial^2=0$ i że homologia
nie zależy od diagramu Heegaarda, potrzebna jest
teoria przestrzeni moduli krzywych pseudoholomorficznych.
Zupełność argumentu analitycznego wymaga między innymi
dopuszczalności diagramu (skończoności sum),
transwersalności, zwartości i sklejania przestrzeni moduli.
Przyjmujemy te wyniki z teorii Floera.
W praktyce obliczeniowej używa się również kombinatorycznych
modeli siatkowych (grid homology), które pozwalają liczyć
niezmienniki Floera dla konkretnych diagramów.
\end{uwaga}

\begin{twierdzenie}[Ozsváth--Szabó; wykrywanie rodzaju]
Niech $K\subset S^3$ będzie węzłem. Wtedy
\[
g(K)=\max\{s\in\mathbb Z:\widehat{HFK}(K,s)\ne0\}.
\]
W szczególności najwyższy niezerowy stopień Alexandera
w homologii Floera wyznacza minimalny rodzaj powierzchni
Seiferta. Jest to wynik zewnętrzny z pracy
P.~Ozsvátha i Z.~Szabó, \emph{Holomorphic disks and genus bounds},
Geom. Topol. 8 (2004), 311--334 (arXiv:math/0311496).
\end{twierdzenie}

\begin{twierdzenie}[Ni; wykrywanie włóknistości]
Niech $K\subset S^3$ będzie węzłem rodzaju $g(K)$.
Węzeł $K$ jest włóknisty wtedy i tylko wtedy, gdy
$\dim_{\mathbb F}\widehat{HFK}(K,g(K))=1$,
gdzie sumujemy wszystkie stopnie Maslova przy ustalonym
stopniu Alexandera. Jest to kryterium włóknistości
z pracy Y.~Ni, \emph{Knot Floer homology detects fibred knots},
Invent. Math. 170 (2007), 577--608 (arXiv:math/0607156).
\end{twierdzenie}

\begin{definicja}[Niezmiennik $\tau$]
Dla węzła $K\subset S^3$ niech $\mathcal F_s$
oznacza podkompleks pełnego $\widehat{CF}(S^3)$
generowany przez elementy o stopniu Alexandera
nie większym niż $s$. Definiujemy
\[
 \tau(K)=\min\{s\in\mathbb Z:
 H_*(\mathcal F_s)\longrightarrow
 \widehat{HF}(S^3;\mathbb F)\cong\mathbb F
 \text{ jest niezerowy}\},
\]
gdzie strzałkę indukuje włączenie.
Minimum istnieje: dla dostatecznie małego $s$
podkompleks jest zerowy, a dla dostatecznie
dużego jest całym kompleksem.
Jedyny niezerowy stopień Maslova grupy docelowej
to zero w standardowej normalizacji.
Niezmienniczość filtracji i definicji przyjmujemy
z pracy P.~Ozsvátha i Z.~Szabó,
\emph{Knot Floer homology and the four-ball genus},
Geom. Topol. 7 (2003), 615--639.
\end{definicja}

\begin{stwierdzenie}[Cykl realizujący niezmiennik $\tau$]
\label{prop:kt-tau-cycle}
Niech $C=\widehat{CF}(S^3)$ z filtracją
$\mathcal F_s$ wyznaczoną przez $K$.
Wtedy $\tau(K)$ jest najmniejszym poziomem filtracji
cyklu reprezentującego niezerową klasę w $H_0(C)$.
Ponadto $\widehat{HFK}(K,\tau(K))$ ma niezerowy
składnik w stopniu Maslova zero.
\end{stwierdzenie}

\begin{proof}
Niech $s=\tau(K)$. Ponieważ włączenie
$\mathcal F_s\subset C$ zachowuje stopień Maslova
i $H_*(C)$ jest jednowymiarowe w stopniu zero,
istnieje cykl $z\in(\mathcal F_s)_0$,
którego klasa w $H_0(C)$ jest niezerowa.
Na mocy minimalności $s$ żaden taki cykl
nie należy do $\mathcal F_{s-1}$.
Niech $\overline z$ oznacza obraz $z$ w
$\mathcal F_s/\mathcal F_{s-1}$.
Jest to cykl, gdyż iloraz jest odwzorowaniem
kompleksów. Gdyby był brzegiem, istniałby
$y\in(\mathcal F_s)_1$ taki, że
$z-\partial y\in\mathcal F_{s-1}$.
Element $z-\partial y$ jest cyklem stopnia zero
reprezentującym w $C$ tę samą niezerową klasę co $z$.
Przeczy to minimalności $s$.
Zatem $[\overline z]\ne0$ w
$H_0(\mathcal F_s/\mathcal F_{s-1})$, który
jest składnikiem $\widehat{HFK}(K,s)$
stopnia Maslova zero.
Odwrotnie każdy cykl z $\mathcal F_j$,
który reprezentuje niezerową klasę w $H_0(C)$,
daje niezerowe odwzorowanie w definicji $\tau$,
więc $j\geq s$. To kończy dowód.
\end{proof}

\begin{uwaga}
Nie wystarczy znaleźć dowolnej niezerowej klasy
w $\widehat{HFK}(K,j)$: jej reprezentant w
ilorazie $\mathcal F_j/\mathcal F_{j-1}$ może nie
podnosić się do cyklu w $C$ albo może reprezentować
zerową klasę w $H_*(C)$. Dlatego definicja $\tau$
wymaga pełnego filtrowanego kompleksu.
\end{uwaga}

\begin{twierdzenie}[Ozsváth--Szabó; własności $\tau$]
Dla dowolnych zorientowanych węzłów $K,K_1,K_2\subset S^3$
zachodzą:
$\tau(U)=0$,
$\tau(K_1\#K_2)=\tau(K_1)+\tau(K_2)$,
$\tau(\overline K)=-\tau(K)$ oraz
$|\tau(K)|\leq g_4^{\mathrm{sm}}(K)$.
Dla dodatnich węzłów torusowych $T(p,q)$
o względnie pierwszych $p,q\geq2$,
$\tau(T(p,q))=(p-1)(q-1)/2$.
\end{twierdzenie}

\begin{przyklad}[Rachunek, który odróżnia TOP od DIFF]
Niech $T=T(2,3)$ będzie dodatnim trójlistnikiem,
a $D=\operatorname{Wh}^{+}(T)$ jego dodatnim,
nieskręconym podwojeniem Whiteheada. Wzorzec
Whiteheada ma liczbę nawinięć zero, a jego
satelita węzła trywialnego jest węzłem trywialnym.
Ze wzoru satelitarnego
$\Delta_{P(K)}(t)\doteq\Delta_{P(U)}(t)\Delta_K(t^w)$
dla liczby nawinięć $w$ i normalizacji $\Delta_K(1)=1$
otrzymujemy $\Delta_D(t)=1$.
Twierdzenie Freedmana o węzłach o trywialnym
wielomianie Alexandera zapewnia lokalnie płasko
zanurzony dysk w $D^4$, a zatem
$g_4^{\mathrm{top}}(D)=0$.
Hedden udowodnił, że dla dodatniego podwojenia
Whiteheada z $t$ skrętami zachodzi
$\tau(D_+(K,t))=1$ dla $t<2\tau(K)$
i zero w przeciwnym przypadku
(\emph{Knot Floer homology of Whitehead doubles},
Geom. Topol. 11 (2007), 2277--2338,
arXiv:math/0606094). Tutaj $\tau(T)=1$ i
$t=0<2$, więc $\tau(D)=1$, a zatem
$1\leq g_4^{\mathrm{sm}}(D)$. Standardową powierzchnię Seiferta dla podwojenia
Whiteheada otrzymujemy z pierścienia o rdzeniu $T$
przez dołączenie wstęgi realizującej zaczep
Whiteheada i łączącej dwa składniki brzegu.
Otrzymana powierzchnia jest zorientowana,
spójna, ma jeden składnik brzegu i
charakterystykę Eulera $0-1=-1$, czyli rodzaj jeden.
Stąd
$g_4^{\mathrm{sm}}(D)\leq1$.
Ostatecznie
$g_4^{\mathrm{top}}(D)=0$ oraz $g_4^{\mathrm{sm}}(D)=1$.
To przykład konkretnego węzła, dla którego kategorie
topologiczna i gładka dają różne odpowiedzi.
\end{przyklad}

\section{Związki z algebrami operatorów i fizyką kwantową}
\label{sec:kt-operator}

Wielomian Jonesa pojawił się w teorii podczynników algebr
von Neumanna, gdzie bada się inkluzje czynników
$N\subset M$ oraz indeks $[M:N]$. Dla podczynnika skończonego indeksu $d=[M:N]$
projekcje Jonesa $p_i$ spełniają
$p_i^2=p_i$, $p_i p_{i\pm1}p_i=d^{-1}p_i$
i komutują dla $|i-j|\ge2$.
Po przeskalowaniu $e_i=\sqrt d\,p_i$ dostajemy
$e_i^2=\sqrt d\,e_i$ i
$e_i e_{i\pm1}e_i=d^{3/2}d^{-1}p_i=e_i$,
czyli relacje algebry Temperleya--Lieba
z parametrem $\delta=\sqrt d$. Jest to kontekst
operatorowy, odrębny od formalnego wyboru
parametru Laurenta poniżej. Z ich pomocą konstruuje się
reprezentacje grup warkoczowych, a funkcjonały śladowe
po odpowiedniej normalizacji dają niezmienniki domknięć
warkoczy.

\begin{definicja}[Algebra Temperleya--Lieba]
Niech $R$ będzie przemiennym pierścieniem z jedynką,
a $\delta\in R$. \emph{Algebrą Temperleya--Lieba}
$TL_n(\delta)$ nazywamy unitalną $R$-algebrę generowaną
przez $1,e_1,\dots,e_{n-1}$, z relacjami
$e_i^2=\delta e_i$,
$e_ie_{i\pm1}e_i=e_i$ dla dopuszczalnych indeksów
i $e_ie_j=e_je_i$ dla $|i-j|\geq2$.
\end{definicja}

\begin{stwierdzenie}[Reprezentacja warkoczy]
Niech $R=\mathbb Z[A,A^{-1}]$ i $\delta=-A^2-A^{-2}$.
Wtedy elementy $b_i:=A\,1+A^{-1}e_i\in TL_n(\delta)$
są odwracalne i określają homomorfizm
$\rho\colon B_n\to TL_n(\delta)^\times$ przez
$\rho(\sigma_i)=b_i$.
\end{stwierdzenie}

\begin{proof}
Ponieważ $e_i^2=\delta e_i$, mamy
$(A1+A^{-1}e_i)(A^{-1}1+Ae_i)
=1+(A^2+A^{-2}+\delta)e_i=1$.
Iloczyn w odwrotnej kolejności ma ten sam wynik,
ponieważ oba czynniki są wielomianami w $e_i$.
Dla $|i-j|\ge2$ relacja $b_ib_j=b_jb_i$
wynika z komutacji $e_i,e_j$.
Relację $b_ib_{i+1}b_i=b_{i+1}b_ib_{i+1}$
sprawdzamy w następującym rachunku.
Są to wszystkie relacje definiujące grupę
warkoczową w prezentacji Artina.
\end{proof}

\begin{przyklad}[Weryfikacja relacji Artina w $TL_3$]
Niech $e=e_i$, $f=e_{i+1}$, $a=A$, $b=A^{-1}$ oraz
$\delta=-a^2-b^2$. Po rozwinięciu iloczynów
$(a+be)(a+bf)(a+be)$ i $(a+bf)(a+be)(a+bf)$
pierwszy iloczyn po redukcji ma postać
\[
a^3 1+(2a^2b+ab^2\delta+b^3)e
+a^2bf+ab^2(ef+fe).
\]
Drugi otrzymujemy przez zamianę $e$ z $f$.
Ich różnica wynosi więc
\[
\bigl(a^2b+ab^2\delta+b^3\bigr)(e-f).
\]
Ponieważ $ab=1$, współczynnik ten wynosi
$a+b\delta+b^3=a+b(-a^2-b^2)+b^3=0$.
Wyrazy zawierające $efe$ i $fef$ redukują się za pomocą
relacji Temperleya--Lieba. Oba iloczyny są zatem równe.
Jest to bezpośrednie rachunkowe uzasadnienie homomorfizmu
$B_n\to TL_n(\delta)^\times$.
\end{przyklad}

\begin{uwaga}[Chern--Simons i anyony]
W kwantowej teorii pola Chern--Simonsa wartości średnie
operatorów Wilsona związanych z zamkniętymi krzywymi
prowadzą do kwantowych niezmienników splotów. W modelach
topologicznych obliczeń kwantowych warkocze trajektorii
nieabelowych anyonów działają poprzez reprezentacje
grup warkoczowych na przestrzeni stanów.
Nie każda formalna reprezentacja algebry
Temperleya--Lieba jest reprezentacją unitarną
na przestrzeni Hilberta. W modelach fizycznych
potrzebny jest odpowiedni wybór parametru
i reprezentacji. Nie należy utożsamiać matematycznego
istnienia reprezentacji z fizyczną realizacją
niezawodnego komputera kwantowego: ta ostatnia wymaga dodatkowych
warunków eksperymentalnych i kontroli dekoherencji.
\end{uwaga}

\section{Biologia, chemia i dynamika płynów}
\label{sec:kt-zastosowania}

\begin{przyklad}[Koliste DNA]
Zamknięta cząsteczka dwuniciowego DNA może tworzyć
węzły i sploty. Enzymy topoizomerazy zmieniają jej
topologię, dokonując kontrolowanego przecięcia,
przełożenia i ponownego połączenia nici. W doświadczeniach
z rekombinacją DNA rozkład typów węzłów produktów
może stanowić informację o mechanizmie działania enzymu.
Ważne zastrzeżenie: rzeczywista cząsteczka ma grubość,
energię sprężystą i może ulegać lokalnym zmianom topologii;
model węzła dotyczy konfiguracji bez przecięć.
\end{przyklad}

\begin{przyklad}[Helikalność pola wirowości]
Niech $\Omega\subset\mathbb R^3$ będzie obszarem,
a $u\colon\Omega\to\mathbb R^3$ gładkim polem prędkości
nieściśliwego płynu, dla którego
$\nabla\cdot u=0$. Wirowością nazywamy
$\omega:=\nabla\times u$.
\emph{Helikalność} ma postać
$\mathcal H(u):=\int_\Omega u(x)\cdot\omega(x)\,
\mathrm d\lambda_3(x)$, jeśli całka istnieje.
Gdy wirowość koncentruje się w rozłącznych, cienkich
rurkach wirowych o strumieniach $\Phi_i$ i osiach
$K_i$, wkład wzajemnego splecenia w przybliżeniu
ma postać $2\sum_{i<j}\Phi_i\Phi_j\operatorname{lk}(K_i,K_j)$.
Jest to wzór w modelu cienkich zamkniętych rurek,
na przykład w $\mathbb R^3$ przy odpowiednim
zanikaniu pola w nieskończoności. Współczynnik
dwa wynika z sumowania obu uporządkowanych par
$(i,j)$ i $(j,i)$ i symetrii liczby splecenia.
Występują również wkłady własne rurek, zależne
od geometrii osi i skręcenia linii wirowości;
nie wyznacza ich sama klasa węzła osi.
Por. H.~K. Moffatt, R.~L. Ricca,
\emph{Helicity and the Călugăreanu invariant},
Proc. R. Soc. Lond. A 439 (1992), 411--429. Pokazuje to, że helikalność jest
geometrycznym odpowiednikiem informacji o spleceniu
linii pola, przy odpowiednich warunkach brzegowych.
\end{przyklad}

\begin{uwaga}[Orbity okresowe pól wektorowych]
Jeśli $X$ jest gładkim polem wektorowym na orientowanej
trójrozmaitości $M$, a $\gamma\colon S^1\hookrightarrow M$
jest prostą zamkniętą orbitą okresową, to jej obraz
jest węzłem w $M$. Typ węzła stanowi niezmiennik
geometryczny orbity względem izotopii układu.
W dynamice hiperbolicznej opisuje się niekiedy całe rodziny
orbit za pomocą rozgałęzionych powierzchni zwanych
\emph{szablonami} (templates), co prowadzi do badań
węzłów realizowanych przez przepływ.
\end{uwaga}

\section{Rozpoznawanie węzłów i algorytmy}
\label{sec:kt-algorytmy}

Teoria węzłów łączy głęboką topologię z obliczeniami:
konkretne pytania o równoważność węzłów można kodować
skończonymi diagramami. Niezmienniki, takie jak kolorowania
Foxa, wyznaczniki i wielomiany, są skutecznymi testami
\emph{nierównoważności}: jeśli dwa węzły mają różne wartości,
na pewno nie są równoważne. Zwykle nie są jednak pełnymi
testami równoważności.

\begin{stwierdzenie}
Jeżeli dla dwóch węzłów $K_1,K_2\subset S^3$ zachodzi
$\Delta_{K_1}(t)\not\doteq\Delta_{K_2}(t)$, to węzły
nie są równoważne. Analogiczny wniosek zachodzi,
gdy $V_{K_1}\ne V_{K_2}$ lub
$\operatorname{Kh}(K_1)\not\cong\operatorname{Kh}(K_2)$
jako odpowiednio dwugradowane grupy.
\end{stwierdzenie}

\begin{proof}
Wszystkie wymienione obiekty są niezmiennikami węzła.
Równoważność wymuszałaby równość wartości niezmienników,
więc ich różność wyklucza równoważność.
\end{proof}

\begin{uwaga}
Problem rozpoznawania węzła trywialnego jest rozstrzygalny
algorytmicznie metodami teorii powierzchni normalnych
w trójrozmaitościach. Istnieją także algorytmy do porównywania
topologii uzupełnień węzłów, a w przypadku hiperbolicznym
istotną rolę odgrywają triangulacje, struktury geometryczne
i obliczenia na grupach fundamentalnych.
Dla praktyki ważne jest rozróżnienie między istnieniem
algorytmu a jego efektywnością na dużych diagramach.
\end{uwaga}



